% GENERATED FILE. Edit canonical lessons, then run npm run generate:books.
% canonical-source-hash: fnv1a64:7f4148584d3bffc5
% canonical-lessons: RU-C369-ekologiya, RU-C369-zagryaznenie, RU-C369-zagryaznyat, RU-C369-pererabatyvat, RU-C369-berech

\chapter{The environment}
\label{ch:ru-the-environment}
\hlchaptermodality{369}

\begin{chapteropening}
\textbf{By the end of this chapter:} \emph{I can talk about the environment, pollution and recycling, and say that we must look after nature.}

Complete the last lesson of chapter 369: I can talk about the environment, pollution and recycling, and say that we must look after nature.
\end{chapteropening}

\section[\texorpdfstring{\ru{экология} (ekológiya)}{экология (ekológiya)}]{\ru{экология} (ekológiya) — ecology, the environment}

\label{lesson:RU-C369-ekologiya}



\begin{quote}
Before the new one: say the Russian for to donate, then the Russian for to need.
\end{quote}

\subsection*{You'll want to know: \ru{экология}}
\textbf{\ru{экология}} (\emph{ekológiya}) — \textquotedblleft{}ecology, the environment\textquotedblright{}.

\begin{quote}
\ru{Мой} \ru{брат} \ru{изучает} \ru{экологию}.
\end{quote}

\begin{quote}
\emph{Moy brat izucháyet ekológiyu.}
\end{quote}

\begin{quote}
My brother studies ecology.
\end{quote}

Feminine. In everyday speech, the state of the environment: \emph{\ru{плохая} \ru{экология}}.

\begin{cousinweb}
\textbf{\ru{экология}} is usually traced to German \emph{Ökologie}, a learned coinage from Greek \emph{oîkos} \textquotedblleft{}a house, a household\textquotedblright{} and \emph{lógos} \textquotedblleft{}word, study\textquotedblright{}: the study of the home we live in.
\end{cousinweb}

\begin{grammarlens}[title={What you've built}]
One more word for telling, explaining and describing at B1.
\end{grammarlens}

\subsection*{Your turn}
\begin{itemize}
\raggedright
  \item \emph{Say it:} \emph{ekológiya}
  \item \emph{Say it:} \emph{ekológiya}, once more
  \item \emph{Say it:} a sentence of your own with \emph{ekológiya}
  \item \emph{Recall:} say the Russian for to donate, then the Russian for to need, then say \emph{ekológiya} again
\end{itemize}

\begin{tcolorbox}[breakable,colback=teal!4,colframe=teal!35!black,title={Before you move on}]
What does \ru{экология} mean? (\textquotedblleft{}Ecology, the environment\textquotedblright{}.) Say it once more.
\end{tcolorbox}
\section[\texorpdfstring{\ru{загрязнение} (zagryaznéniye)}{загрязнение (zagryaznéniye)}]{\ru{загрязнение} (zagryaznéniye) — pollution}

\label{lesson:RU-C369-zagryaznenie}



\begin{quote}
Before the new one: say the Russian for to need, then the Russian for ecology.
\end{quote}

\subsection*{You'll want to know: \ru{загрязнение}}
\textbf{\ru{загрязнение}} (\emph{zagryaznéniye}) — \textquotedblleft{}pollution\textquotedblright{}.

\begin{quote}
\ru{Загрязнение} \ru{воздуха} — \ru{серьёзная} \ru{проблема}.
\end{quote}

\begin{quote}
\emph{Zagryaznéniye vózdukha — seryóznaya probléma.}
\end{quote}

\begin{quote}
Air pollution is a serious problem.
\end{quote}

Neuter. \emph{\ru{Загрязнение} \ru{воды}} — water pollution.

\begin{cousinweb}
\textbf{\ru{загрязнение}} is built on \emph{\ru{загрязнить}} \textquotedblleft{}to pollute\textquotedblright{}, from \emph{\ru{за}-} and \emph{\ru{грязь}} \textquotedblleft{}dirt, mud\textquotedblright{}, the root of \emph{\ru{грязный}} \textquotedblleft{}dirty\textquotedblright{}.
\end{cousinweb}

\begin{grammarlens}[title={What you've built}]
One more word for telling, explaining and describing at B1.
\end{grammarlens}

\subsection*{Your turn}
\begin{itemize}
\raggedright
  \item \emph{Say it:} \emph{zagryaznéniye}
  \item \emph{Say it:} \emph{zagryaznéniye}, once more
  \item \emph{Say it:} a sentence of your own with \emph{zagryaznéniye}
  \item \emph{Recall:} say the Russian for to need, then the Russian for ecology, then say \emph{zagryaznéniye} again
\end{itemize}

\begin{tcolorbox}[breakable,colback=teal!4,colframe=teal!35!black,title={Before you move on}]
What does \ru{загрязнение} mean? (\textquotedblleft{}Pollution\textquotedblright{}.) Say it once more.
\end{tcolorbox}
\section[\texorpdfstring{\ru{загрязнять} (zagryaznyát)}{загрязнять (zagryaznyát)}]{\ru{загрязнять} (zagryaznyát) — to pollute}

\label{lesson:RU-C369-zagryaznyat}



\begin{quote}
Before the new one: say the Russian for ecology, then the Russian for pollution.
\end{quote}

\subsection*{You'll want to know: \ru{загрязнять}}
\textbf{\ru{загрязнять}} (\emph{zagryaznyát}) — \textquotedblleft{}to pollute\textquotedblright{}.

\begin{quote}
\ru{Машины} \ru{загрязняют} \ru{воздух}.
\end{quote}

\begin{quote}
\emph{Mashíny zagryaznyáyut vózdukh.}
\end{quote}

\begin{quote}
Cars pollute the air.
\end{quote}

Imperfective; perfective \emph{\ru{загрязнить}}. It takes the accusative.

\begin{cousinweb}
\textbf{\ru{загрязнять}} is built from \emph{\ru{за}-} and \emph{\ru{грязь}} \textquotedblleft{}dirt, mud\textquotedblright{}, the root of \emph{\ru{грязный}}: to make a thing dirty.
\end{cousinweb}

\begin{grammarlens}[title={What you've built}]
One more word for telling, explaining and describing at B1.
\end{grammarlens}

\subsection*{Your turn}
\begin{itemize}
\raggedright
  \item \emph{Say it:} \emph{zagryaznyát}
  \item \emph{Say it:} \emph{zagryaznyát}, once more
  \item \emph{Say it:} a sentence of your own with \emph{zagryaznyát}
  \item \emph{Recall:} say the Russian for ecology, then the Russian for pollution, then say \emph{zagryaznyát} again
\end{itemize}

\begin{tcolorbox}[breakable,colback=teal!4,colframe=teal!35!black,title={Before you move on}]
What does \ru{загрязнять} mean? (\textquotedblleft{}To pollute\textquotedblright{}.) Say it once more.
\end{tcolorbox}
\section[\texorpdfstring{\ru{перерабатывать} (pererabátyvat)}{перерабатывать (pererabátyvat)}]{\ru{перерабатывать} (pererabátyvat) — to recycle, to process}

\label{lesson:RU-C369-pererabatyvat}



\begin{quote}
Before the new one: say the Russian for pollution, then the Russian for to pollute.
\end{quote}

\subsection*{You'll want to know: \ru{перерабатывать}}
\textbf{\ru{перерабатывать}} (\emph{pererabátyvat}) — \textquotedblleft{}to recycle, to process\textquotedblright{}.

\begin{quote}
\ru{В} \ru{нашем} \ru{городе} \ru{перерабатывают} \ru{пластик}.
\end{quote}

\begin{quote}
\emph{V náshem górode pererabátyvayut plástik.}
\end{quote}

\begin{quote}
In our town they recycle plastic.
\end{quote}

Imperfective; perfective \emph{\ru{переработать}}. It also means to work overtime.

\begin{cousinweb}
\textbf{\ru{перерабатывать}} is built from \emph{\ru{пере}-} \textquotedblleft{}over again\textquotedblright{} and \emph{\ru{работать}} \textquotedblleft{}to work\textquotedblright{}, from \emph{\ru{работа}}, Proto-Slavic \emph{*orbota}: to work a material over into something new.
\end{cousinweb}

\begin{grammarlens}[title={What you've built}]
One more word for telling, explaining and describing at B1.
\end{grammarlens}

\subsection*{Your turn}
\begin{itemize}
\raggedright
  \item \emph{Say it:} \emph{pererabátyvat}
  \item \emph{Say it:} \emph{pererabátyvat}, once more
  \item \emph{Say it:} a sentence of your own with \emph{pererabátyvat}
  \item \emph{Recall:} say the Russian for pollution, then the Russian for to pollute, then say \emph{pererabátyvat} again
\end{itemize}

\begin{tcolorbox}[breakable,colback=teal!4,colframe=teal!35!black,title={Before you move on}]
What does \ru{перерабатывать} mean? (\textquotedblleft{}To recycle, to process\textquotedblright{}.) Say it once more.
\end{tcolorbox}
\section[\texorpdfstring{\ru{беречь} (beréch)}{беречь (beréch)}]{\ru{беречь} (beréch) — to look after, to protect; to save}

\label{lesson:RU-C369-berech}



\begin{quote}
Before the new one: say the Russian for to pollute, then the Russian for to recycle.
\end{quote}

\subsection*{You'll want to know: \ru{беречь}}
\textbf{\ru{беречь}} (\emph{beréch}) — \textquotedblleft{}to look after, to protect; to save\textquotedblright{}.

\begin{quote}
\ru{Нужно} \ru{беречь} \ru{природу}.
\end{quote}

\begin{quote}
\emph{Núzhno beréch priródu.}
\end{quote}

\begin{quote}
We must look after nature.
\end{quote}

Imperfective; present \emph{\ru{берегу}, \ru{бережёшь}, \ru{берегут}}; past \emph{\ru{берёг}, \ru{берегла}}. \emph{\ru{Беречь} \ru{себя}} — to take care of oneself.

\begin{cousinweb}
\textbf{\ru{беречь}} is inherited from Proto-Slavic \emph{*berťi}, from PIE \emph{*b\textsuperscript{h}erg\textsuperscript{h}-} \textquotedblleft{}to defend, to protect\textquotedblright{} (EtymDB).
\end{cousinweb}

\begin{grammarlens}[title={What you've built}]
One more word for telling, explaining and describing at B1.
\end{grammarlens}

\subsection*{Your turn}
\begin{itemize}
\raggedright
  \item \emph{Say it:} \emph{beréch}
  \item \emph{Say it:} \emph{beréch}, once more
  \item \emph{Say it:} a sentence of your own with \emph{beréch}
  \item \emph{Recall:} say the Russian for to pollute, then the Russian for to recycle, then say \emph{beréch} again
\end{itemize}

\begin{tcolorbox}[breakable,colback=teal!4,colframe=teal!35!black,title={Before you move on}]
What does \ru{беречь} mean? (\textquotedblleft{}To look after, to protect; to save\textquotedblright{}.) Say it once more.
\end{tcolorbox}
